\section{Experimental Methodology}
\label{sec:method}

Let $\mathcal{E}_T$ be the set of episodes in a run. We report attack
success rate (ASR) and detection rate (DR):
\begin{align}
  \mathrm{ASR}
  &=
  \frac{\bigl|\{\,e\in\mathcal{E}_T : \text{exfiltration succeeded}\,\}\bigr|}
       {\bigl|\{\,e\in\mathcal{E}_T : \text{attack attempted}\,\}\bigr|},
  \label{eq:asr}\\[2pt]
  \mathrm{DR}
  &=
  \frac{\bigl|\{\,e\in\mathcal{E}_T : \text{at least one detection}\,\}\bigr|}
       {\bigl|\{\,e\in\mathcal{E}_T : \text{attack attempted}\,\}\bigr|}.
  \label{eq:dr}
\end{align}
Secondary quantities include resource costs, recovery time, coalition and
deception rates, strategy diversity and behavioural novelty. Co-evolutionary
pressure averages the change in attacker and defender signatures between
consecutive episodes,
\begin{equation}
  \mathrm{CEP}
  \;=\;
  \frac{1}{T-1}
  \sum_{t=1}^{T-1}
  \frac{\bigl\|s_A^{(t+1)}-s_A^{(t)}\bigr\|_2
        +\bigl\|s_D^{(t+1)}-s_D^{(t)}\bigr\|_2}{2},
  \label{eq:cep}
\end{equation}
and emergent security risk (ESR) compares failure under interaction to
failure under isolation,
\begin{equation}
  \mathrm{ESR}
  \;=\;
  \frac{P(\mathrm{failure}\mid \mathrm{interaction})}
       {P(\mathrm{failure}\mid \mathrm{isolation})}.
  \label{eq:esr}
\end{equation}
Values above $1$ indicate interaction-induced risk. ESR requires paired
runs and is reserved for future population studies (H3).

Baselines follow the B1--B5 ladder: rule-based agents (B1/B2), then
language-model or \emph{hybrid} agents (LLM proposals constrained by
legal capability hints, with heuristic takeover on stall) at static,
episodic and persistent memory levels. Heuristic adaptive policies with
stealth/vigilance parameters validate the loop without a language model.
Table~\ref{tab:matrix} lists the full E1--E8 programme; this paper reports
the one-versus-one ladder and B3--B5 cells. We test seven hypotheses (H1--H7)
on adaptation gain, defender recovery, interaction risk, memory regime,
emergence, population effects and non-monotonic model capability.

Each cell is replicated over multiple random seeds; we report means and
95\% confidence intervals (Student-$t$ on seed means). Primary cells use
six seeds; supplementary ablations use two or three seeds and are read
directionally. Measurements are collected at three levels: step
trajectories, episode records (outcome, costs, strategy signature) and
run aggregates (ASR, DR, CEP, dynamics class). Table~\ref{tab:episode-record}
illustrates one episode row; Fig.~\ref{fig:episode-dynamics} plots the
same persistent run through time.

\begin{table}[!t]
\caption{Experimental matrix. A\,=\,attacker, D\,=\,defender.}
\label{tab:matrix}
\centering
\setlength{\tabcolsep}{3pt}
\begin{tabular}{@{}clccc@{}}
\toprule
\textbf{ID} & \textbf{Agents} & \textbf{Adapt.} & \textbf{Comm.} & \textbf{Goal} \\
\midrule
E1 & $1$A+$1$D & none       & none    & baseline \\
E2 & $1$A+$1$D & episodic   & none    & adaptation \\
E3 & $1$A+$1$D & episodic   & direct  & interaction \\
E4 & $3$A+$3$D & episodic   & direct  & population \\
E5 & $5$A+$5$D & episodic   & graph   & scaling \\
E6 & $5$A+$5$D & persistent & graph   & co-evolution \\
E7 & $5$A+$5$D & persistent & graph   & coalition \\
E8 & $10$A+$10$D & persistent & dynamic & emergence \\
\bottomrule
\end{tabular}
\end{table}

\begin{table}[!t]
\caption{Illustrative per-episode record (persistent heuristic run, episode~20).}
\label{tab:episode-record}
\centering
\setlength{\tabcolsep}{3pt}
\begin{tabular}{@{}llr@{}}
\toprule
\textbf{Quantity} & \textbf{Interpretation} & \textbf{Value} \\
\midrule
Attack success & Exfiltration achieved & no \\
Detected & At least one alert & no \\
Attacker / defender cost & Resource use & $21.6$ / $24.7$ \\
Attacker stealth & Signature parameter & $0.70$ \\
Defender vigilance & Signature parameter & $0.19$ \\
\bottomrule
\end{tabular}
\end{table}

All configurations, analysis scripts and aggregate results are available
in the public repository cited in Appendix~\ref{app:repro}.
